\documentclass[11pt, a4paper]{article}

\usepackage[a4paper, top=2.5cm, bottom=2.5cm, left=2cm, right=2cm]{geometry}
\usepackage{amsmath}
\usepackage{amssymb}
\usepackage{amsfonts}
\usepackage{mathtools}
\usepackage{booktabs}
\usepackage{enumitem}
\usepackage[colorlinks=true, linkcolor=blue, citecolor=blue, urlcolor=blue]{hyperref}

\title{\textbf{Information-Topological Theory: Minimal Foundational Formulation}}
\author{\textbf{Damyan Damyanov}\\
\small\textit{Independent Researcher, Hisarya, Bulgaria}\\
\small\href{mailto:damione1@gmail.com}{damione1@gmail.com}}
\date{\today}

\begin{document}

\maketitle

\section*{Abstract}
We present the minimal foundational formulation of Information-Topological Theory (ITT). By deliberately isolating standard mathematical invariants from ITT-specific physical postulates, we establish the topological level selection $k_{\min} = 6$ purely through a $\mathbb{Z}_6$-equivariant descent condition on the configuration space $L(6,1) = S^3/\mathbb{Z}_6$, without relying on phenomenological Standard Model parameter fits.

\newpage

\section{Information-Topological Theory: Minimal Foundational Formulation}

\subsection{Scope and Logical Status}
The Information-Topological Theory (ITT) is formulated as a timeless canonical gauge theory on a three-dimensional configuration manifold. The fundamental configuration space is the gauge-orbit space
\begin{equation}
\mathcal C_{\rm ITT} = \mathcal A(S^3)/\mathcal G(S^3),
\end{equation}
where $\mathcal A(S^3)$ denotes the space of smooth $SU(2)$ connections on $S^3$ and $\mathcal G(S^3)$ the corresponding gauge transformation group.
The theory deliberately separates standard mathematical results from ITT-specific postulates. In particular, the integer quantization of the Chern--Simons level is standard; the selection of the minimal level $k=6$ is an ITT equivariant-descent condition.

\subsection{I. Timeless Canonical Configuration Space}
The fundamental canonical variables are
\begin{equation}
A_i^a(x), \qquad E_a^i(x),
\end{equation}
with canonical Poisson brackets
\begin{equation}
\{A_i^a(x), E_b^j(y)\} = \delta_i^j \delta_b^a \delta^{(3)}(x-y),
\end{equation}
and
\begin{equation}
\{A_i^a, A_j^b\} = 0, \qquad \{E_a^i, E_b^j\} = 0.
\end{equation}
No external physical time is postulated at the fundamental level. A parameter $\tau$ may be introduced only as a bookkeeping parameter for curves in configuration space.

\subsection{II. Chern--Simons Functional and ITT Constraint}
Let
\begin{equation}
F = dA + A \wedge A
\end{equation}
and
\begin{equation}
B_a^i = \frac{1}{2}\epsilon^{ijk}F_{jk,a}.
\end{equation}
The $SU(2)$ generators are normalized by
\begin{equation}
T_a = \frac{\sigma_a}{2i}, \qquad \operatorname{tr}(T_a T_b) = -\frac{1}{2}\delta_{ab}.
\end{equation}
Define the Chern--Simons functional
\begin{equation}
S_{\rm CS}^{(0)}[A] = \int_{S^3} \operatorname{tr} \left( A \wedge dA + \frac{2}{3} A \wedge A \wedge A \right),
\end{equation}
with functional derivative
\begin{equation}
\frac{\delta S_{\rm CS}^{(0)}}{\delta A_i^a} = -B_a^i.
\end{equation}
The fundamental ITT constraint is
\begin{equation}
\boxed{
\mathcal C_a^i \equiv E_a^i + \frac{k}{4\pi}B_a^i = E_a^i - \frac{k}{4\pi} \frac{\delta S_{\rm CS}^{(0)}}{\delta A_i^a} \approx 0.
}
\end{equation}
The corresponding first-order timeless action is
\begin{equation}
S_{\rm ITT} = \int d\tau \int_{S^3} d^3x \left[ E_a^i \frac{dA_i^a}{d\tau} - N_i^a \mathcal C_a^i - A_0^a \mathcal G_a \right].
\end{equation}

\subsection{III. Constraint Algebra}
Because $\mathcal C$ is a canonical momentum shift by an exact functional one-form,
\begin{equation}
\boxed{
\{\mathcal C_a^i(x), \mathcal C_b^j(y)\} = 0
}
\end{equation}
on the closed manifold $S^3$.
The magnetic field satisfies the non-Abelian Bianchi identity
\begin{equation}
D_i B_a^i \equiv 0.
\end{equation}
Consequently,
\begin{equation}
D_i \mathcal C_a^i = D_i E_a^i \equiv \mathcal G_a.
\end{equation}
Thus the Gauss constraint is a differential consequence of the ITT constraint:
\begin{equation}
\boxed{
\mathcal G_a = D_i \mathcal C_a^i.
}
\end{equation}
At the classical level the constraint system is therefore first-class with a differential reducibility relation. At the quantum level the corresponding identity is understood with a fixed gauge-covariant operator ordering.

\subsection{IV. Formal Exact Chern--Simons State}
Canonical quantization is defined by
\begin{equation}
\widehat E_a^i = -i\frac{\delta}{\delta A_i^a}.
\end{equation}
The formal wave functional
\begin{equation}
\boxed{
\Psi_k[A] = \mathcal N \exp\left[ i\frac{k}{4\pi} S_{\rm CS}^{(0)}[A] \right]
}
\end{equation}
satisfies
\begin{equation}
\widehat{\mathcal C}_a^i \Psi_k[A] = 0
\end{equation}
because
\begin{equation}
\widehat E_a^i \Psi_k = -\frac{k}{4\pi} B_a^i \Psi_k.
\end{equation}
Consequently, for the positive quadratic constraint
\begin{equation}
\widehat{\mathcal H}_{\rm ITT} = \frac{1}{2} \widehat{\mathcal C}_a^i \widehat{\mathcal C}_i^a,
\end{equation}
with a specified regularization and operator ordering,
\begin{equation}
\boxed{
\widehat{\mathcal H}_{\rm ITT}\Psi_k = 0.
}
\end{equation}
This is a formal exact solution of the ITT canonical constraint. A complete Hilbert-space construction, regularization and self-adjoint realization are separate analytical problems.

\subsection{V. Unit Topological Sector}
For the $SU(2)$ BPST instanton on $D_R^4 \subset \mathbb R^4$, with scale $\rho$, define
\begin{equation}
c = \frac{R^2}{R^2+\rho^2}.
\end{equation}
With the above trace convention and orientation,
\begin{equation}
Q(R) = -\frac{1}{8\pi^2} \int_{D_R^4} \operatorname{tr}(F \wedge F) = 3c^2 - 2c^3 = \frac{R^4(R^2+3\rho^2)}{(R^2+\rho^2)^3}.
\end{equation}
For
\begin{equation}
A|_{S^3} = c \, g^{-1}dg,
\end{equation}
the corresponding Chern--Simons transgression gives the same quantity,
\begin{equation}
\frac{1}{2\pi} S_{\rm CS}(A|_{S^3}) = 3c^2 - 2c^3.
\end{equation}
In the asymptotic limit $R/\rho \rightarrow \infty$, one obtains the unit topological sector
\begin{equation}
\boxed{Q \rightarrow 1.}
\end{equation}
The corresponding APS cylinder representation yields the net spectral-flow index
\begin{equation}
\boxed{
\operatorname{SF}(\mathcal D_{A_c}) = \operatorname{ind}\left(\mathcal D_{[0,1]\times S^3}^{\rm APS}\right) = 1,
}
\end{equation}
with $c$ interpreted solely as a parameter of a family of three-dimensional gauge configurations and not as physical time.

\subsection{VI. Standard Level Quantization}
Under a large $SU(2)$ gauge transformation with winding number
\begin{equation}
w \in \pi_3(SU(2)) \cong \mathbb Z,
\end{equation}
the normalized Chern--Simons action shifts by an integer multiple of $2\pi$. Therefore gauge invariance of the quantum amplitude requires
\begin{equation}
\boxed{
k \in \mathbb Z.
}
\end{equation}
This is a standard Chern--Simons quantization condition and does not by itself select $k=6$.

\subsection{VII. Lens-Space Topological Sector}
The ITT quotient geometry is defined by the free diagonal action
\begin{equation}
(z_1,z_2) \mapsto \left( e^{2\pi i/6}z_1, e^{2\pi i/6}z_2 \right)
\end{equation}
on $S^3 \subset \mathbb C^2$. The resulting lens space is
\begin{equation}
\boxed{
L(6,1) = S^3/\mathbb Z_6,
}
\end{equation}
with $\pi_1(L(6,1)) \cong \mathbb Z_6$. Let $\rho_1 : \mathbb Z_6 \rightarrow SU(2)$ be the fundamental non-trivial flat representation,
\begin{equation}
\rho_1(\gamma) = \begin{pmatrix} e^{2\pi i/6} & 0 \\ 0 & e^{-2\pi i/6} \end{pmatrix}.
\end{equation}
With the fixed orientation and trace convention, its Chern--Simons invariant is
\begin{equation}
\boxed{
\operatorname{cs}(\rho_1) = -\frac{1}{6} \pmod{\mathbb Z}.
}
\end{equation}
The value follows from the rational intersection/linking form associated with the $(-6)$-framed surgery description of $L(6,1)$.

\subsection{VIII. ITT Equivariant Descent Condition}
The ITT quantum state is required to admit a $\mathbb Z_6$-equivariant lift to its prequantum line bundle. Denote the generator of $\mathbb Z_6$ by $\zeta$.

The ITT equivariant structure is required to satisfy
\begin{equation}
\boxed{
\chi_k(\zeta) = \exp\left( 2\pi i k \operatorname{cs}(\rho_1) \right).
}
\end{equation}
Using $\operatorname{cs}(\rho_1) = -\frac{1}{6}$, the fiber character becomes
\begin{equation}
\boxed{
\chi_k(\zeta) = e^{-2\pi i k/6}.
}
\end{equation}
This equation is the defining ITT equivariant-lift condition. It is not asserted as a universal theorem of ordinary $SU(2)$ Chern--Simons theory.

\subsection{IX. ITT Level Selection Rule}
The ITT physical state is required to descend to the quotient configuration space with trivial stabilizer action. Therefore
\begin{equation}
\chi_k(\zeta) = 1.
\end{equation}
Hence
\begin{equation}
e^{-2\pi i k/6} = 1
\end{equation}
and therefore
\begin{equation}
\boxed{
k \in 6\mathbb Z.
}
\end{equation}
The minimal positive ITT level is
\begin{equation}
\boxed{
k_{\min} = 6.
}
\end{equation}
Thus the level selection is not derived from a phenomenological relation such as $3 \times 2$. It is an ITT-specific equivariant descent condition associated with the chosen $\mathbb Z_6$ quotient geometry.

\subsection{X. Independent Soliton Sector}
The number of particle generations is not used in the derivation of $k$. Instead, ITT treats
\begin{equation}
N_{\rm gen} = 3
\end{equation}
as an independent dynamical prediction of the soliton sector.
The future objective is to construct an ITT energy functional
\begin{equation}
E_{\rm ITT}[n,A]
\end{equation}
whose stable topological sectors satisfy
\begin{equation}
Q_H = 1, 2, 3
\end{equation}
while higher-charge configurations are dynamically unstable or energetically decomposable.
Until such a result is established, $N_{\rm gen} = 3$ remains an ITT postulate or target prediction and is not used in the proof of $k_{\min} = 6$.

\subsection{XI. Foundational Logical Chain}
The minimal ITT foundation is summarized by the sequence:
\begin{equation}
\begin{array}{c}
\mathcal C_{\rm ITT} = \mathcal A(S^3)/\mathcal G(S^3) \\
\Downarrow \\
\boxed{ \mathcal C_a^i = E_a^i + \frac{k}{4\pi}B_a^i, \qquad \{\mathcal C, \mathcal C\} = 0 } \\
\Downarrow \\
\boxed{ D_i\mathcal C_a^i = \mathcal G_a } \\
\Downarrow \\
\boxed{ \Psi_k[A] = \mathcal N e^{\,i k S_{\rm CS}^{(0)}[A]/4\pi} } \\
\Downarrow \\
\boxed{ L(6,1) = S^3/\mathbb Z_6, \qquad \operatorname{cs}(\rho_1) = -\frac{1}{6} } \\
\Downarrow \\
\boxed{ \chi_k(\zeta) = e^{-2\pi i k/6} } \\
\Downarrow \\
\boxed{ \chi_k(\zeta) = 1 \Longrightarrow k \in 6\mathbb Z } \\
\Downarrow \\
\boxed{ k_{\min} = 6. }
\end{array}
\end{equation}

The independent soliton problem begins after this foundational chain:
\begin{equation}
E_{\rm ITT} \longrightarrow \text{stable Hopf sectors } (Q_H = 1,2,3) \longrightarrow N_{\rm gen} = 3.
\end{equation}

\subsection{Foundational Status}
The minimal theory therefore distinguishes three levels of statement:
\begin{enumerate}
\item \textbf{Standard mathematics:} gauge theory on $S^3$, BPST topology, APS spectral flow, Chern--Simons integer level quantization, lens-space topology, and rational Chern--Simons invariants.
\item \textbf{ITT defining structure:} the helicity constraint, the timeless canonical formulation, the $\mathbb Z_6$-equivariant quantum structure, and the requirement of trivial equivariant descent.
\item \textbf{Future ITT predictions:} the dynamical derivation of $N_{\rm gen} = 3$, the particle mass spectrum, coupling constants, and independent experimental observables.
\end{enumerate}
The foundational claim of ITT is therefore:
\begin{equation}
\boxed{ \text{timeless gauge configuration space} + \text{topological constraint} + \text{$\mathbb Z_6$ equivariant descent} \Longrightarrow k_{\min} = 6. }
\end{equation}
No Standard-Model parameter fit and no phenomenological factor $3 \times 2$ is used in this foundational selection of the Chern--Simons level.

\end{document}